\documentclass[12pt]{article}
\usepackage[T1]{fontenc}
\usepackage[utf8]{inputenc}
\usepackage{lmodern}
\usepackage{textcomp}
\usepackage{enumitem}
\usepackage{geometry}
\usepackage{graphicx}
\usepackage{amsmath, amssymb}
\geometry{margin=1in}
\begin{document}
\begin{center}
Cybersecurity - Graduate Level Assessment 10 \
Total Marks: 40 \
Answer all questions.
\end{center}

\section*{Multiple Choice (10 marks)}
\begin{enumerate}[label=\arabic*.]
\item To verify a digital signature we need the
\begin{enumerate}[label=(\alph*)]
    \item Sender's Private key
    \item Sender's Public key
    \item Receiver's Private key
    \item Receiver's Public key
\end{enumerate}
\item In a \_\_\_\_\_\_\_\_\_\_\_\_\_ attack, the extra data that holds some specific instructions in the memory for actions is projected by a cyber-criminal or penetration tester to crack the system.
\begin{enumerate}[label=(\alph*)]
    \item Phishing
    \item MiTM
    \item Buffer-overflow
    \item Clickjacking
\end{enumerate}
\item The \_\_\_\_\_\_\_\_\_\_\_\_\_\_ is categorized as an unknown segment of the Deep Web which has been purposely kept hidden \& is inaccessible using standard web browsers.
\begin{enumerate}[label=(\alph*)]
    \item Haunted web
    \item World Wide Web
    \item Dark web
    \item Surface web
\end{enumerate}
\item Which of the following does authorization aim to accomplish?
\begin{enumerate}[label=(\alph*)]
    \item Restrict what operations/data the user can access
    \item Determine if the user is an attacker
    \item Flag the user if he/she misbehaves
    \item Determine who the user is
\end{enumerate}
\item \_\_\_\_\_\_\_\_\_\_\_\_\_\_\_\_\_\_\_\_ is the anticipation of unauthorized access or break to computers or data by means of wireless networks.
\begin{enumerate}[label=(\alph*)]
    \item Wireless access
    \item Wireless security
    \item Wired Security
    \item Wired device apps
\end{enumerate}
\end{enumerate}

\section*{True / False (10 marks)}
\textit{Answer True or False}
\begin{enumerate}[label=\arabic*.]
\setcounter{enumi}{5}
\item True or False: Memory leakage is a coding mistake that is primarily caused by inadequate processing power in applications.
\item True or False: The Metasploit framework simplifies the exploitation of vulnerabilities through a user-friendly, point-and-click interface.
\item True or False: The one-time pad key can be computed using the formula k = m xor c.
\item True or False: Nmap classifies port states as active, inactive, or standby when assessing network connections.
\item True or False: DaCryptic is widely recognized as a remote Trojan similar to Troya in functionality.
\end{enumerate}

\section*{Long Form Response (20 marks)}
\textit{Answer each question with a clear and complete explanation.}
\begin{enumerate}[label=\arabic*.]
\setcounter{enumi}{10}
\item Discuss the four-step authentication process that occurs between a wireless user and an access point, commonly known as the 4-way handshake, and analyze its significance in cybersecurity.
\item Discuss the principles of mutation-based fuzzing in cybersecurity, focusing on how it generates new inputs by modifying existing ones and its implications for software testing and vulnerability discovery.
\end{enumerate}

\vfill
\noindent\textit{End of Paper}
\end{document}